\documentclass[border=6pt]{standalone}

\usepackage{array}
\usepackage[table]{xcolor}
\usepackage{amsmath}

\renewcommand{\arraystretch}{1.28}
\setlength{\arrayrulewidth}{0.5pt}

\newcommand{\casprompt}{\raisebox{0.15ex}{\rule{0.8ex}{0.8ex}}\hspace{0.55em}}
\newcommand{\casin}[1]{{\ttfamily\detokenize{#1}}}
\newcommand{\casout}[1]{\(\displaystyle #1\)}

\begin{document}

\begin{tabular}{
  |>{\raggedright\arraybackslash}p{4.5cm}
  |>{\raggedright\arraybackslash}p{4.5cm}|
}
\hline
\rowcolor{gray!10}
\casprompt \casin{f(x)=4x-1} & \casout{f(x)=4x-1} \\
\hline

\rowcolor{gray!10}
\casprompt \casin{g(x)=x^3-6x^2+15x-7} & \casout{g(x)=x^3-6x^2+15x-7} \\
\hline

\rowcolor{gray!10}
\casprompt \casin{h(x)=g(x)-f(x)} & \casout{h(x)=x^3-6x^2+11x-6} \\
\hline

\rowcolor{gray!10}
\casprompt \casin{factor(h(x))} & \casout{(x-3)\cdot(x-2)\cdot(x-1)} \\
\hline

\rowcolor{gray!10}
\casprompt \casin{j(x)=f(x)+3*h(x)} & \casout{j(x)=3x^3-18x^2+37x-19} \\
\hline

\rowcolor{gray!10}
\casprompt \casin{j(1)} & \casout{3} \\
\hline

\rowcolor{gray!10}
\casprompt \casin{j(2)} & \casout{7} \\
\hline

\rowcolor{gray!10}
\casprompt \casin{j(3)} & \casout{11} \\
\hline

\rowcolor{gray!10}
\casprompt & \\
\hline
\end{tabular}

\end{document}